\documentclass[12pt,a4paper]{article}

\usepackage[portuguese]{babel}
\usepackage[utf8]{inputenc}
\usepackage[T1]{fontenc}
\usepackage{amsmath}
\usepackage{amssymb} 
\usepackage{braket}
\usepackage{hyperref}
\setlength{\parskip}{0.8em}
\setlength{\parindent}{0pt}
\begin{document}
\section{Introdução}

Os sistemas operativos modernos distribuem carga entre núcleos de CPU recorrendo a heurísticas clássicas
que, apesar de eficazes em produção, tomam decisões subótimas em cenários de elevada contenção de recursos.

Em paralelo, métodos de otimização quântica, nomeadamente o Quantum Approximate Optimization Algorithm (QAOA)~\cite{Farhi2014QAOA},
têm sido investigados como abordagens alternativas para problemas combinatoriais formuláveis
como QUBO (\textit{Quadratic Unconstrained Binary Optimization})~\cite{Kochenberger2014QUBO,Glover2019QUBO}.

No entanto, a era NISQ (\textit{Noisy Intermediate-Scale Quantum}) impõe limitações de escala e ruído que
tornam difícil avaliar quando um escalonador quântico oferece vantagens concretas sobre alternativas
clássicas~\cite{Preskill2018}. Preskill antecipa que esta barreira será ultrapassada com hardware da ordem
do \textit{megaquop} --- máquinas capazes de executar $10^6$ operações quânticas com erro
tolerável~\cite{Preskill2025} --- fornecendo um referencial concreto para situar a relevância futura
deste tipo de abordagem.

Esta tese descreve o desenvolvimento de uma plataforma de testes híbrida quantum-clássica para o problema
de atribuição de processos a núcleos de CPU. A plataforma não substitui o escalonador do kernel; isola a
decisão de atribuição processo-núcleo, modela-a como um QUBO, resolve-a com QAOA via PennyLane~\cite{Bergholm2018PennyLane}, e valida
os resultados contra um \textit{solver} exato de \textit{brute force} e simulated annealing~\cite{Kirkpatrick1983SA}. O objetivo não é demonstrar
superioridade quântica, mas caracterizar o comportamento do QAOA neste domínio em função da escala do problema.

\subsection{Formulação do Problema}

Dado um conjunto de processos com utilização de CPU conhecida e um número fixo de núcleos, o sistema deve
atribuir cada processo a exatamente um núcleo. O objetivo de otimização é o equilíbrio de carga: a carga
total de cada núcleo deve aproximar-se o mais possível da carga média global. A solução deve ainda respeitar
restrições de atribuição exclusiva --- nenhum processo pode ficar por atribuir ou ser atribuído a múltiplos
núcleos simultaneamente.

Expressar estas restrições e este objetivo numa formulação QUBO, compatível com otimizadores quânticos,
constitui o problema central desta tese.

% ─────────────────────────────────────────────
\section{Formulação QUBO para Escalonamento}

\subsection{CPU Scheduling em Sistemas Operativos}
Um escalonador de CPU determina que processo ou \textit{thread} executa em cada núcleo do processador
em cada instante. Escalonadores como o Completely Fair Scheduler (CFS) do Linux, consideram
simultaneamente latência, \textit{throughput}, prioridade, afinidade de núcleo e localização de cache,
tomando decisões em microssegundos com base em heurísticas clássicas~\cite{LinuxCFS}.

No entanto, a decisão de atribuição de processos a núcleos é fundamentalmente um problema de partição combinatorial. Dado um conjunto de $N$ processos com cargas $w_i$ e $K$ núcleos, o número de atribuições possíveis cresce como $K^N$, tornando a procura exaustiva para instâncias reais.

Esta tese adota uma formulação simplificada deste problema: cada \textit{workload entity} possui
uma utilização de CPU conhecida, e o objetivo é distribuí-las pelos núcleos disponíveis minimizando o
desequilíbrio de carga. Esta abstração isola a componente combinatorial do escalonamento sem replicar a complexidade de um escalonador de kernel completo.

\subsection{Formulação QUBO}

Um problema QUBO (\textit{Quadratic Unconstrained Binary Optimization}) ~\cite{Kochenberger2014QUBO,Glover2019QUBO} consiste em minimizar
uma expressão quadrática sobre variáveis binárias $x_i \in \{0, 1\}$, representada de forma compacta como:

\begin{equation}
    \min_{\mathbf{x} \in \{0,1\}^n} \mathbf{x}^\top Q\, \mathbf{x} = \sum_{i} \sum_{j} Q_{ij} \cdot x_i \cdot x_j
\end{equation}

onde $Q \in \mathbb{R}^{n \times n}$ é uma matriz cujos coeficientes codificam o problema a resolver.
Expandindo o somatório, surgem dois tipos de termos:

\begin{itemize}
    \item \textbf{Termos diagonais} ($i = j$): como $x_i^2 = x_i$ para variáveis binárias,
    $Q_{ii} \cdot x_i \cdot x_i = Q_{ii} \cdot x_i$. Estes termos associam um custo individual
    a cada variável --- um coeficiente $Q_{ii} < 0$ incentiva o otimizador quântico a escolher $x_i = 1$; um coeficiente
    $Q_{ii} > 0$ incentiva $x_i = 0$.

    \item \textbf{Termos fora da diagonal} ($i \neq j$): $Q_{ij} \cdot x_i \cdot x_j$ contribui
    para a energia apenas quando $x_i = x_j = 1$. Um coeficiente positivo penaliza a ativação
    simultânea de $x_i$ e $x_j$; um coeficiente negativo recompensa-a.
\end{itemize}

A energia de uma solução $\mathbf{x}$ é o valor $\mathbf{x}^\top Q\, \mathbf{x}$.
A solução ótima $\mathbf{x}^*$ é a atribuição binária que minimiza esta energia no espaço de
$2^n$ configurações possíveis.

\subsection{Construção da Matriz Q}

Dado um conjunto de $N$ processos com cargas $w_i$ e $K$ núcleos, define-se uma
variável binária $x_{ij} \in \{0,1\}$ para cada par processo-núcleo. Esta variável
toma o valor $1$ quando o processo $i$ é atribuído ao núcleo $j$, e o valor $0$
caso contrário.
Na implementação, a variável $x_{ij}$ é armazenada no índice linear $iK+j$ da matriz $Q$.

O objetivo original do problema é equilibrar a carga entre núcleos. Para isso,
mede-se o desequilíbrio através do desvio quadrático total das cargas relativamente
à carga média:

\begin{equation}
    \min \sum_{j=0}^{K-1}(L_j - L_{avg})^2.
\end{equation}

onde
\begin{equation}
    L_{avg} = \frac{W_{total}}{K}.
\end{equation}

Para obter a formulação QUBO, este termo é expandido e combinado com penalidades que impõem
atribuições válidas. O resultado é codificado na matriz $Q$, de modo que o \textit{solver}
minimize a energia
\begin{equation}
    E(\mathbf{x}) = \mathbf{x}^{\top}Q\mathbf{x}.
\end{equation}

Assim, minimizar a energia QUBO corresponde, por construção, a procurar uma atribuição
simultaneamente equilibrada e válida.

Para cada núcleo $j$, a carga total é dada por:
\begin{equation}
    L_j = \phi_j + \sum_{i=0}^{N-1} w_i x_{ij},
\end{equation}
onde $\phi_j$ representa carga já comprometida nesse núcleo, relevante para a secção da
decomposição iterativa. Na primeira iteração, ou num QUBO sem carga fixa, tem-se
$\phi_j = 0$.

Substituindo $L_j$ no termo de um núcleo e usando $x_{ij}^2=x_{ij}$, obtém-se:
\begin{equation}
\begin{aligned}
(L_j - L_{avg})^2
&= \left(\phi_j + \sum_{i=0}^{N-1}w_ix_{ij} - L_{avg}\right)^2 \\
&= \underbrace{\sum_{i=0}^{N-1} d_{ij}x_{ij}}_{\text{termos diagonais}}
 + 2\underbrace{\sum_{0 \leq i < k}^{N-1} w_iw_kx_{ij}x_{kj}}_{\text{termos de co-localização}}
 + \underbrace{(\phi_j - L_{avg})^2}_{\text{constante}},
\end{aligned}
\end{equation}
onde
\begin{equation}
    d_{ij}=w_i^2 + 2w_i\phi_j - 2L_{avg}w_i.
\end{equation}

Estes coeficientes são então
adicionados à matriz $Q$. Os termos diagonais estão associados
à ativação individual de uma variável $x_{ij}$; os termos de co-localização aparecem
quando dois processos distintos são atribuídos ao mesmo núcleo.

Além do objetivo de balanceamento, é necessário garantir que cada processo é atribuído a
exatamente um núcleo. Para isso adiciona-se a penalidade one-hot:
\begin{equation}
    P \sum_{i=0}^{N-1}\left(\sum_{j=0}^{K-1}x_{ij} - 1\right)^2.
\end{equation}
Esta penalidade segue a prática comum de converter restrições de igualdade em termos
quadráticos penalizados em QUBO~\cite{Glover2019QUBO}. Neste caso, contribui com $-P$
na diagonal de cada variável e com $+P$ entre variáveis do mesmo processo em núcleos diferentes.

Combinando o objetivo de balanceamento e a penalidade one-hot, cada entrada da matriz
é determinada pelo tipo de relação entre as duas variáveis consideradas:
\begin{equation}
Q[x_{ij},x_{kl}] =
\begin{cases}
w_i^2 + 2w_i\phi_j - 2L_{avg}w_i - P, & \text{se } i=k \text{ e } j=l,\\
w_iw_k, & \text{se } i \neq k \text{ e } j=l,\\
P, & \text{se } i=k \text{ e } j \neq l,\\
0, & \text{se } i \neq k \text{ e } j \neq l.
\end{cases}
\end{equation}
A primeira linha corresponde à contribuição individual de cada variável. A segunda
linha penaliza processos distintos que partilham o mesmo núcleo, enquanto a terceira
corresponde à penalidade one-hot entre núcleos diferentes do mesmo processo. A última
linha indica que variáveis de processos e núcleos diferentes não interagem diretamente.

\subsection{Exemplo Ilustrativo}

Considere-se $N=2$ processos e $K=2$ núcleos, com cargas $w_0=0.3$ e $w_1=0.5$,
sem carga fixa ($\phi_j=0$). As variáveis ficam ordenadas como:
\begin{equation}
\begin{array}{c|cccc}
\text{índice} & 0 & 1 & 2 & 3 \\
\hline
\text{variável} & x_{00} & x_{01} & x_{10} & x_{11}
\end{array}
\end{equation}
A carga média é $L_{avg}=(0.3+0.5)/2=0.4$.

Antes de adicionar a penalidade one-hot, os termos diagonais do balanceamento são:
\begin{align}
Q_{0,0}^{bal}=Q_{1,1}^{bal}
&= w_0^2 - 2L_{avg}w_0
= 0.09 - 2(0.4)(0.3)
= -0.15,\\
Q_{2,2}^{bal}=Q_{3,3}^{bal}
&= w_1^2 - 2L_{avg}w_1
= 0.25 - 2(0.4)(0.5)
= -0.15.
\end{align}

Para o núcleo $0$, a expansão de $(L_0 - L_{avg})^2$ contém o termo de
co-localização $2w_0w_1x_{00}x_{10}=0.30x_{00}x_{10}$; de forma análoga, a expansão de $(L_1 - L_{avg})^2$ contém
$2w_0w_1x_{01}x_{11}=0.30x_{01}x_{11}$.

Assim, a matriz do objetivo de balanceamento é:
\begin{equation}
Q_{bal} =
\begin{pmatrix}
-0.15 & 0 & 0.15 & 0 \\
0 & -0.15 & 0 & 0.15 \\
0.15 & 0 & -0.15 & 0 \\
0 & 0.15 & 0 & -0.15
\end{pmatrix}.
\end{equation}

Com $P=1.0$, a penalidade one-hot subtrai $1.0$ a cada termo diagonal e adiciona $1.0$
entre as variáveis do mesmo processo em núcleos distintos: $x_{00}$ com $x_{01}$, e
$x_{10}$ com $x_{11}$. A matriz final fica:
\begin{equation}
Q =
\begin{pmatrix}
-1.15 & 1.0 & 0.15 & 0 \\
1.0 & -1.15 & 0 & 0.15 \\
0.15 & 0 & -1.15 & 1.0 \\
0 & 0.15 & 1.0 & -1.15
\end{pmatrix}.
\end{equation}

Para verificar o efeito dos termos de co-localização, comparem-se duas soluções viáveis:
A, com os processos distribuídos ($x_{00}=x_{11}=1$), e B, com ambos os processos no
núcleo $0$ ($x_{00}=x_{10}=1$). Usando $E(\mathbf{x})=\mathbf{x}^{\top}Q\mathbf{x}$:
\begin{align}
E_A &= Q_{0,0} + Q_{3,3}
= -1.15 + (-1.15)
= \mathbf{-2.30},\\
E_B &= Q_{0,0} + Q_{2,2} + Q_{0,2} + Q_{2,0}
= -1.15 + (-1.15) + 0.15 + 0.15
= \mathbf{-2.00}.
\end{align}

Como $E_A < E_B$, a solução distribuída tem menor energia. A matriz penaliza, portanto,
a concentração dos dois processos no mesmo núcleo, favorecendo uma distribuição mais
equilibrada.

Depois de construída a matriz $Q$, o problema de escalonamento passa a estar expresso
numa forma compatível com otimizadores QUBO. O próximo passo será resolver esta matriz, isto é, procurar a bitstring $\mathbf{x}$ que minimiza a energia $E(\mathbf{x})=\mathbf{x}^{\top}Q\mathbf{x}$. Neste trabalho, essa etapa é realizada
com QAOA~\cite{Farhi2014QAOA}, implementado em PennyLane~\cite{Bergholm2018PennyLane}.

\section{Resolução do QUBO com QAOA}
O QAOA é um algoritmo variacional híbrido para otimização combinatória, no qual
um circuito quântico é combinado com um otimizador clássico para
minimizar o valor esperado de um Hamiltoniano de custo~\cite{Farhi2014QAOA,Cerezo2021VQA}.
Neste projeto, a implementação do circuito e da otimização é feita com PennyLane~\cite{Bergholm2018PennyLane}.

Importa realçar que a execução foi realizada em simulação, e não em hardware
quântico real. Esta opção permite evitar a latência associada aos dispositivos quânticos remotos, bem como os efeitos de ruído físico. Na implementação, o solver tenta
usar o backend \texttt{lightning.gpu} do PennyLane quando disponível, recorrendo a
\texttt{lightning.qubit} como alternativa quando a execução em GPU não é possível.

A abordagem implementada atribui um qubit lógico a cada par processo-núcleo. Assim,
uma instância com $N$ processos e $K$ núcleos requer $N \cdot K$ qubits.

\subsection{Conversão do QUBO em Hamiltoniano}
O QAOA trabalha com um Hamiltoniano de custo $H_C$ e não diretamente com a matriz $Q$. Assim, a energia QUBO

\begin{equation}
    E(\mathbf{x}) = \mathbf{x}^{\top}Q\mathbf{x}
\end{equation}

tem de ser convertida em forma de operadores quânticos. Esta conversão segue a
relação usual entre variáveis binárias e variáveis de \textit{spin} usada em formulações

Ising de problemas combinatórios~\cite{Lucas2014Ising}. Para isso, cada variável binária $x_i \in \{0,1\}$ é mapeada para operadores de Pauli $Z_i$ usando a substituição:

\begin{equation}
    x_i = \frac{I - Z_i}{2}.
\end{equation}

Esta substituição permite representar os valores binários $0$ e $1$ através dos
dois estados da base computacional, $\lvert 0\rangle$ e $\lvert 1\rangle$. Quando o qubit está no estado $\lvert 0\rangle$, o operador $Z_i$ tem valor
$+1$, resultando em $x_i = 0$; quando está no estado $\lvert 1\rangle$, o operador $Z_i$
tem valor $-1$, resultando em $x_i = 1$.

Assim, um termo linear da diagonal do QUBO transforma-se em:

\begin{equation}
    Q_{ii}x_i
    =
    Q_{ii}\frac{I - Z_i}{2}
    =
    \frac{Q_{ii}}{2}I - \frac{Q_{ii}}{2}Z_i.
\end{equation}

Para os termos quadráticos, considerando $q_{ij}=Q_{ij}+Q_{ji}$, obtém-se:

\begin{equation}
    q_{ij}x_ix_j
    =
    q_{ij}\frac{I - Z_i}{2}\frac{I - Z_j}{2}
    =
    \frac{q_{ij}}{4}
    \left(I - Z_i - Z_j + Z_iZ_j\right).
\end{equation}

Desta forma, o Hamiltoniano de custo pode ser escrito como:

\begin{equation}
    H_C =
    cI
    +
    \sum_i h_i Z_i
    +
    \sum_{i<j} J_{ij} Z_iZ_j,
\end{equation}

onde $c$ é uma constante, $h_i$ são os coeficientes lineares associados a cada
qubit, e $J_{ij}$ representa a interação entre pares de qubits.

Uma vez obtido o Hamiltoniano de custo $H_C$ este passa a definir a energia associada a cada estado. No QAOA, esta informação é usada através de um,a operação unitária de custo:

\begin{equation}
    U_C(\gamma) = e^{-i\gamma H_C},
\end{equation}

onde $\gamma$ é um parâmetro variacional. Esta operação pode ser vista como uma evoluação temporal sob o Hamiltoniano $H_C$ durante um tempo controlado por $\gamma$. Como $H_C$ é diagonal na base computacional, cada bitstring recebe uma fase dependente da sua energia:

\begin{equation}
    U_C(\gamma)\lvert x\rangle
    =
    e^{-i\gamma E(x)}\lvert x\rangle.
\end{equation}

Logo a camada de custo não altera diretamente a probabilidade de cada bitstring, ela altera as suas fases relativas. Essas fases tornam-se utéis quando é aplicada a camada de mistura, responsável por distribuir amplitude entre soluções.

A camada de mistura é definida por outro Hamiltoniano $H_M$ e pela unidade

\begin{equation}
\label{eq_mixer_op}
    U_M(\beta) = e^{-i\beta H_M},
\end{equation}

onde $\beta$ é também um parâmetro variacional. Enquanto $U_C$ traduz informação
sobre a qualidade das soluções através das fases, $U_M$ permite explorar novas
soluções ao misturar amplitudes entre estados da base computacional.

\subsection{Mixer X e XY}
O Hamiltoniano $H_M$ presente em \eqref{eq_mixer_op} é chamado de \textit{mixer} e temos disponivel dois tipos: X e XY.

A diferença entre eles essencialmente afeta se o QAOA explora todo o espaço binário ou apenas o espaço de atribuições válidas.

\subsubsection{Mixer X}
Este é o padrão do QAOA. O Hamiltoniano de forma simplificada é:

\begin{equation}
    H_M^X=\sum_i X_i
\end{equation}

onde $X_i$ é o operador de Pauli X aplicado ao qubit $i$.

Este mixer define que o estado inicial deve ser a superposição uniforme sobre todos os bitstrings:

\begin{equation}
    \ket{s}=\frac{1}{\sqrt{2^n}} \sum_x \ket{x}
\end{equation}

No entanto, este espaço contém muitos estado inválidos. Por exemplo, com $N=1$ e $K=2$ os estados possiveis são:
\begin{equation}
    00, 01, 10, 11
\end{equation}

No entanto só dois são válidos:
\begin{equation}
    01, 10
\end{equation}

E por essa razão que para o nosso caso de uso, o \textit{mixer XY} é o mais relevante.

\subsubsection{Mixer XY}

Em vez de operar no espaço total de soluções, ele opera somente nos qubits que representam os possiveis núcleos desse processo.

Para cada processo $i$, há um bloco de $K$ qubits:
\begin{equation}
    x_{i0},x_{i1},...,x_{i,K-1},
\end{equation}

O mixer XY atua somente entre pares de qubits desse mesmo bloco através da operação:

\begin{equation}
   X_{ij}X_{ik} + Y_{ij}Y_{ik}
\end{equation}

que permite trocar amplitude entre estados como:
\begin{equation}
   \ket{10} <-> \ket{01}
\end{equation}

mas sem criar os estados inválidos
\begin{equation}
   \ket{00} <-> \ket{11}
\end{equation}

Dito isto, concluimos que o mixer XY é o mais válido para este projeto, transformando um espaço de soluções de tamanho 

$$
2^{N \cdot K}
$$

num espaço de tamanho

$$
K^N
$$


\bibliographystyle{plain}
\bibliography{references}

\end{document}
